% SPDX-License-Identifier: CC-BY-SA-4.0
% Part: Syntactic analysis
% Section: Decidability of algebraic languages
% Exercise: Revision exercise

\begingroup

\begin{exercice}[Révisions -- Analyseur LR]\label{exo:parsing/decidability/revision}

  Pour chacune des grammaires suivantes :
  \begin{itemize}
  \item Dessiner l'automate fini sous-jacent obtenu lors de la conception de l'analyseur LR.
  \item Dessiner l'automate à pile de l'analyseur LR.
  \item Si l'automate à pile possède des conflits, dessiner l'automate à pile de l'analyseur SLR.
  \end{itemize}

  \begin{question}
  \item $\left\langle \{a, b\}, \{S, A, B\}, S,
    \left\{\begin{array}{lll}
    S   &\rightarrow& A \$\\
    A   &\rightarrow& Aa \mid  B\\
    B   &\rightarrow& Bb \mid  \varepsilon\\
    \end{array}\right\}\right\rangle$

  \item $\left\langle \{ +, (, ), 1\}, \{A, B\}, A,
    \left\{\begin{array}{lll}
    A   &\rightarrow& B + A \mid B \\
    B   &\rightarrow& 1
    \end{array}\right\}\right\rangle$

  \item $\left\langle \{ +, (, ), 1\}, \{A, B\}, A,
    \left\{\begin{array}{lll}
    A   &\rightarrow& A + B \mid B \\
    B   &\rightarrow& (A) \mid 1
    \end{array}\right\}\right\rangle$

  \item $\left\langle \{ +, *, (, ), 1\}, \{A, B, C\}, A,
    \left\{\begin{array}{lll}
    A   &\rightarrow& A + B \mid  B \\
    B   &\rightarrow& B * C \mid C \\
    C   &\rightarrow& (A) \mid 1
    \end{array}\right\}\right\rangle$

    
  \end{question}

\end{exercice}

\begin{correction}

  \begin{question}
    
  \item
    \begin{itemize}
    \item Automate sous-jacent :

      \begin{tikzpicture}[automaton, baseline=(0.base), x=40mm, y=20mm]
        \state[initial, rectangle] (0) at (1,2) {
          $0=\left\{\begin{array}{@{}r@{}c@{}l@{}}
          S &\rightarrow&\bullet A\$\\
          A &\rightarrow&\bullet Aa \mid \bullet B\\
          B &\rightarrow&\bullet Bb \mid \bullet \varepsilon
          \end{array}\right.$};
        
        \state[rectangle] (1) at (1,1) {
          $1=\left\{\begin{array}{@{}r@{}c@{}l@{}}
          S &\rightarrow& A\bullet\$\\
          A &\rightarrow& A\bullet a \mid B\\
          B &\rightarrow& Bb \mid \varepsilon
          \end{array}\right.$};
        
        \state[rectangle] (2) at (2,1) {
          $2=\left\{\begin{array}{@{}r@{}c@{}l@{}}
          S &\rightarrow& A\$\\
          A &\rightarrow& Aa \mid B\bullet\\
          B &\rightarrow& B\bullet b \mid \varepsilon
          \end{array}\right.$};
        
        \state[accepting, rectangle] (3) at (0,1) {
          $3=\left\{\begin{array}{@{}r@{}c@{}l@{}}
          S &\rightarrow& A\$\bullet\\
          A &\rightarrow& Aa \mid B\\
          B &\rightarrow& Bb \mid \varepsilon
          \end{array}\right.$};
        
        \state[rectangle] (4) at (1,0) {
          $4=\left\{\begin{array}{@{}r@{}c@{}l@{}}
          S &\rightarrow& A\$\\
          A &\rightarrow& Aa\bullet \mid B\\
          B &\rightarrow& Bb \mid \varepsilon
          \end{array}\right.$};
        
        \state[rectangle] (5) at (2,0) {
          $5=\left\{\begin{array}{@{}r@{}c@{}l@{}}
          S &\rightarrow& A\$\\
          A &\rightarrow& Aa \mid B\\
          B &\rightarrow& Bb\bullet \mid \varepsilon
          \end{array}\right.$};
        
        \path (0) edge node       {$A$}  (1);
        \path (0) edge node       {$B$}  (2);
        \path (1) edge node       {$\$$} (3);
        \path (1) edge node       {$a$}  (4);
        \path (2) edge node       {$b$}  (5);
      \end{tikzpicture}
      
    \item Automate à pile (un conflit par décallage/réduction dans l'état $2$) :

      \begin{tikzpicture}[pushdown, baseline=(0.base), x=40mm, y=20mm, initial text=$0$]
        \state[initial above, rectangle] (0) at (1,2) {
          $0=\left\{\begin{array}{@{}r@{}c@{}l@{}}
          S &\rightarrow&\bullet A\$\\
          A &\rightarrow&\bullet Aa \mid \bullet B\\
          B &\rightarrow&\bullet Bb \mid \bullet \varepsilon
          \end{array}\right.$};
        
        \state[rectangle] (1) at (1,1) {
          $1=\left\{\begin{array}{@{}r@{}c@{}l@{}}
          S &\rightarrow& A\bullet\$\\
          A &\rightarrow& A\bullet a \mid B\\
          B &\rightarrow& Bb \mid \varepsilon
          \end{array}\right.$};
        
        \state[rectangle] (2) at (2,1) {
          $2=\left\{\begin{array}{@{}r@{}c@{}l@{}}
          S &\rightarrow& A\$\\
          A &\rightarrow& Aa \mid B\bullet\\
          B &\rightarrow& B\bullet b \mid \varepsilon
          \end{array}\right.$};
        
        \state[accepting, rectangle] (3) at (0,1) {
          $3=\left\{\begin{array}{@{}r@{}c@{}l@{}}
          S &\rightarrow& A\$\bullet\\
          A &\rightarrow& Aa \mid B\\
          B &\rightarrow& Bb \mid \varepsilon
          \end{array}\right.$};
        
        \state[rectangle] (4) at (1,0) {
          $4=\left\{\begin{array}{@{}r@{}c@{}l@{}}
          S &\rightarrow& A\$\\
          A &\rightarrow& Aa\bullet \mid B\\
          B &\rightarrow& Bb \mid \varepsilon
          \end{array}\right.$};
        
        \state[rectangle] (5) at (2,0) {
          $5=\left\{\begin{array}{@{}r@{}c@{}l@{}}
          S &\rightarrow& A\$\\
          A &\rightarrow& Aa \mid B\\
          B &\rightarrow& Bb\bullet \mid \varepsilon
          \end{array}\right.$};
        
        \path (1) edge            node {\smPAtrans{\$}{\varepsilon}{3}}   (3);
        \path (1) edge[bend left] node {\smPAtrans{a}{\varepsilon}{4}}    (4);
        \path (2) edge[bend left] node {\smPAtrans{b}{\varepsilon}{5}}    (5);
        \path (0) edge            node {\smPAtrans{\varepsilon}{0}{02}}   (2);
        \path (2) edge            node {\smPAtrans{\varepsilon}{02}{01}}  (1);
        \path (4) edge[bend left] node {\smPAtrans{\varepsilon}{014}{01}} (1);
        \path (5) edge[bend left] node {\smPAtrans{\varepsilon}{025}{02}} (2);

      \end{tikzpicture}
    \end{itemize}

  \item Ici, le conflit par décalage/réduction de l'état $3$ ne peut pas être résolu par une simple fusion de transitions :
    il faut observer le symbole $\$$ sans le consommer, ce qui ne peut pas être représenté directement dans le formalisme
    des automates à pile adopté ici, mais se met naturellement en \oe uvre dans l'implémentation d'un analyseur SLR.

    \begin{tikzpicture}[pushdown, baseline=(0.base), x=50mm, y=20mm, initial text=$0$]
      \state[initial above, rectangle] (0) at (0,3) {
        $0=\left\{\begin{array}{@{}r@{}c@{}l@{}}
        S &\rightarrow&\bullet A\$\\
        A &\rightarrow&\bullet B + A \mid \bullet B\\
        B &\rightarrow&\bullet 1
        \end{array}\right.$};
      
      \state[rectangle] (1) at (0,0) {
        $1=\left\{\begin{array}{@{}r@{}c@{}l@{}}
        S &\rightarrow& A\bullet\$\\
        A &\rightarrow& B + A \mid  B\\
        B &\rightarrow& 1
        \end{array}\right.$};
      
      \state[rectangle] (2) at (0,2) {
        $2=\left\{\begin{array}{@{}r@{}c@{}l@{}}
        S &\rightarrow& A\$\\
        A &\rightarrow& B + A \mid  B\\
        B &\rightarrow& 1\bullet
        \end{array}\right.$};
      
      \state[rectangle] (3) at (0,1) {
        $3=\left\{\begin{array}{@{}r@{}c@{}l@{}}
        S &\rightarrow& A\$\\
        A &\rightarrow& B\bullet + A \mid  B\bullet\\
        B &\rightarrow& 1
        \end{array}\right.$};
      
      \state[accepting, rectangle] (4) at (1,0) {
        $4=\left\{\begin{array}{@{}r@{}c@{}l@{}}
        S &\rightarrow& A\$\bullet\\
        A &\rightarrow& B + A \mid  B\\
        B &\rightarrow& 1
        \end{array}\right.$};
      
      \state[rectangle] (5) at (1,2) {
        $5=\left\{\begin{array}{@{}r@{}c@{}l@{}}
        S &\rightarrow& A\$\\
        A &\rightarrow& \bullet B + \bullet A \mid \bullet B\\
        B &\rightarrow& \bullet 1
        \end{array}\right.$};

      \state[rectangle] (6) at (1,1) {
        $6=\left\{\begin{array}{@{}r@{}c@{}l@{}}
        S &\rightarrow& A\$\\
        A &\rightarrow& B + A \bullet\mid  B\\
        B &\rightarrow& 1
        \end{array}\right.$};

      \path (0) edge             node         {\smPAtrans{1}{\varepsilon}{2}}     (2);
      \path (5) edge             node         {\smPAtrans{1}{\varepsilon}{2}}     (2);
      \path (3) edge             node[sloped] {\smPAtrans{+}{\varepsilon}{5}}     (5);
      \path (1) edge             node         {\smPAtrans{\$}{\varepsilon}{4}}    (4);
      
      \path (2) edge             node         {\smAlign{\smPAtrans{\varepsilon}{02}{03}\smPAtrans{\varepsilon}{52}{53}}}   (3);
      \path (3) edge             node         {\smPAtrans{\varepsilon}{03}{01}}   (1);
      \path (3) edge             node         {\smPAtrans{\varepsilon}{53}{56}}   (6);
      \path (6) edge             node[sloped] {\smPAtrans{\varepsilon}{5356}{56}} (1);
      \path (6) edge[loop right] node         {\smPAtrans{\varepsilon}{0356}{01}} (6);

    \end{tikzpicture}

  \item Même technique...
    
  \item Même technique...

  \end{question}

\end{correction}

\endgroup
\endinput
